\tocsection{\textbf{8/27/26 Lecture 1: }Introduction}

\textbf{Deductive Reasoning:} Drawing conclusions from premises (``axioms''). This is non-ampliative, meaning no new information is created. Using deductive reasoning is ``risk-free'' in the sense that if the premises are true, all conclusions are true as well.\\

\textbf{Inductive Reasoning: } Using observations to make general claims. This is ampliative, meaning that you are creating new information. However, there are ``risks'' to inductive reasoning, in the sense that one's conclusions may still be wrong even if the premises hold. 
\begin{itemize}
    \item \textbf{Problem of Induction (Hume): }How can one justify making an inductive claim? One can argue using the uniformity principle that what happens today will happen tomorrow, but this is justifying inductive reasoning using inductive reasoning! Thus, let's just assume it works. 
\end{itemize}


Now, all of statistics and science is build on inductive reasoning, so how do we even make sure the claims we make are valid (in some sense)? The way to do this is by turning it into a deductive reasoning problem, and there are two main ways to do that. 
\begin{itemize}
    \item \textbf{Idea \#1 Bayesian Reasoning: }Start with prior beliefs and use mathematics and probability theory to update these beliefs.
    \begin{itemize}
        \item \textbf{Pros:} We can actually say something like ``there is a 95\% chance our hypothesis is true''
        \item \textbf{Cons:} We must make an assumption on what the prior beliefs are, which we cannot justify.
    \end{itemize}
    \item \textbf{Idea \#2 Inductive Behavior (Frequentist): }Reason directly from the data, and quantize error.
    \begin{itemize}
        \item \textbf{Pros:} No assumptions needed, relies strictly on data.
        \item \textbf{Cons:} We cannot quantize the probability our hypothesis is true, but instead we can say that a process works $\%95$ of the time under repetition.
    \end{itemize}
\end{itemize}
